\chapter{Pattern Matching\label{chap:PatternMatching}}

The construct \texttt{$e$ IN $p$} tests whether the value calculated by the expression $e$,
which is referred to as the \emph{discriminant},
matches a value specified by the \emph{pattern} $p$.

\ExampleDef{Patterns}
\listingref{Patterns} shows examples of some patterns.
\ASLListing{Examples of patterns}{Patterns}{\definitiontests/Patterns.asl}

\ChapterOutline
\begin{itemize}
  \item \FormalRelationsRef{Patterns} defines the formal relations for patterns;
  \item \secref{MatchingAllValues} defines the match-all-values pattern;
  \item \secref{MatchingASingleValue} defines the match-a-single-value pattern;
  \item \secref{MatchingARangeOfIntegers} defines the match-a-range-of-integers pattern;
  \item \secref{MatchingAnUpperBoundedRangeOfIntegers} defines the match-integers-less-than pattern;
  \item \secref{MatchingALowerBoundedRangeOfIntegers} defines the match-integers-greater-than pattern;
  \item \secref{MatchingABitmask} defines the bitmask pattern;
  \item \secref{MatchingPatternLists} defines positive and negative matching against pattern lists.
\end{itemize}

%%%%%%%%%%%%%%%%%%%%%%%%%%%%%%%%%%%%%%%%%%%%%%%%%%%%%%%%%%%%%%%%%%%%%%%%%%%%%%%%
\FormalRelationsDef{Patterns}
\paragraph{Syntax:} Patterns are grammatically derived from $\Npattern$.

\paragraph{Abstract Syntax:} Patterns are derived in the abstract grammar from $\pattern$
  and generated from $\buildpattern$.
  \hypertarget{build-pattern}{}
  The function
\[
  \buildpattern(\overname{\parsenode{\Npattern}}{\vparsednode}) \;\aslto\;
  \overname{\pattern}{\vastnode}
\]
transforms a pattern parse node $\vparsednode$ into a pattern AST node $\vastnode$.

\paragraph{Typing:} Patterns are annotated by $\annotatepattern$.
\RenderRelation{annotate_pattern}
\paragraph{Semantics:} Patterns are evaluated by $\evalpattern$.
\RenderRelation{eval_pattern}

\section{Matching All Values\label{sec:MatchingAllValues}}
\hypertarget{any-pattern}{}
\ASLListing{Matching any value}{semantics-pall}{\semanticstests/SemanticsRule.PAll.asl}

\subsection{Syntax}
\begin{flalign*}
\Npattern \derives\ & \Tminus &
\end{flalign*}

\subsection{Abstract Syntax}
\RenderTypes[remove_hypertargets]{pattern_all}

\ASTRuleDef{PAll}
\begin{mathpar}
\inferrule{}{
  \buildpattern(\Npattern(\Tminus)) \astarrow
  \overname{\PatternAll}{\vastnode}
}
\end{mathpar}

\subsection{Typing}
\TypingRuleDef{PAll}
The pattern \verb|-| in \listingref{semantics-pall} is well-typed.

\RenderProseAndFormally{annotate_pattern_all}
\CodeSubsection{\PAllBegin}{\PAllEnd}{../Typing.ml}

\subsection{Semantics}
\SemanticsRuleDef{PAll}
\ExampleDef{Evaluation of a Match-all Pattern}
In \listingref{semantics-pall}, \texttt{match\_me} evaluates to \True,
since \texttt{-} matches any value and \texttt{42} in particular.


\RenderProseAndFormally{eval_pattern_PAll}
\CodeSubsection{\EvalPAllBegin}{\EvalPAllEnd}{../Interpreter.ml}

\section{Matching a Single Value\label{sec:MatchingASingleValue}}
\subsection{Syntax}
\begin{flalign*}
\Npattern \derives\ & \Nexpr &
\end{flalign*}

\subsection{Abstract Syntax}
\RenderTypes[remove_hypertargets]{pattern_single}

\ASTRuleDef{PSingle}
\begin{mathpar}
\inferrule{}{
  \buildpattern(\Npattern(\punnode{\Nexpr})) \astarrow
  \overname{\PatternSingle(\astof{\vexpr})}{\vastnode}
}
\end{mathpar}

\subsection{Typing}
\TypingRuleDef{PSingle}
\ExampleDef{Typing Single-expression Patterns}
\listingref{typing-psingle} shows examples of well-typed single-expression patterns.
\ASLListing{Typing single-expression patterns}{typing-psingle}{\typingtests/TypingRule.PSingle.asl}

\listingref{typing-psingle-bad} shows an ill-typed single-expression pattern.
\ASLListing{Ill-typed single-expression pattern}{typing-psingle-bad}{\typingtests/TypingRule.PSingle.bad.asl}


\RenderProseAndFormally{annotate_pattern_single}
\CodeSubsection{\PSingleBegin}{\PSingleEnd}{../Typing.ml}

\subsection{Semantics}
\SemanticsRuleDef{PSingle}
\ExampleDef{Evaluation of a Single-expression Pattern}
In \listingref{semantics-psingle},
\texttt{match\_true} evaluates to \True, since \texttt{42} matches \verb|{42}|, whereas \texttt{match\_false} evaluates to \False, since \texttt{42} does not match \verb|{3}|.
\ASLListing{Matching against a single value}{semantics-psingle}{\semanticstests/SemanticsRule.PSingle.asl}


\RenderProseAndFormally{eval_pattern_PSingle}
\CodeSubsection{\EvalPSingleBegin}{\EvalPSingleEnd}{../Interpreter.ml}

\section{Matching a Range of Integers\label{sec:MatchingARangeOfIntegers}}
\subsection{Syntax}
\begin{flalign*}
\Npattern \derives\ & \Nexpr \parsesep \Tslicing \parsesep \Nexpr &
\end{flalign*}

\subsection{Abstract Syntax}
\RenderTypes[remove_hypertargets]{pattern_range}

\ASTRuleDef{PRange}
\begin{mathpar}
\inferrule{}{
  {
    \begin{array}{r}
  \buildpattern(\Npattern(\namednode{\veone}{\Nexpr}, \Tslicing, \namednode{\vetwo}{\Nexpr})) \astarrow\\
  \overname{\PatternRange(\astof{\veone}, \astof{\vetwo})}{\vastnode}
    \end{array}
  }
}
\end{mathpar}

\subsection{Typing}
\TypingRuleDef{PRange}
\ExampleDef{Typing Range Patterns}
\listingref{typing-prange} shows examples of well-typed range patterns.
\ASLListing{Well-typed range patterns}{typing-prange}{\typingtests/TypingRule.PRange.asl}

\listingref{typing-prange-bad} shows ill-typed range patterns.
\ASLListing{Ill-typed range patterns}{typing-prange-bad}{\typingtests/TypingRule.PRange.bad.asl}


\RenderProseAndFormally{annotate_pattern_range}
\CodeSubsection{\PRangeBegin}{\PRangeEnd}{../Typing.ml}

\subsection{Semantics}
\SemanticsRuleDef{PRange}
\ExampleDef{Evaluation of a Range Pattern}
In \listingref{semantics-prange},
\texttt{match\_true} evaluates to \True, since \texttt{42} is in the range given by
\texttt{3..42},
whereas \texttt{match\_false} evaluates to \False, since \texttt{1} is outside the
range given by \texttt{3..42}.
\ASLListing{Matching against a range of values}{semantics-prange}{\semanticstests/SemanticsRule.PRange.asl}


\RenderProseAndFormally{eval_pattern_PRange}
\CodeSubsection{\EvalPRangeBegin}{\EvalPRangeEnd}{../Interpreter.ml}

\section{Matching an Upper Bounded Range of Integers\label{sec:MatchingAnUpperBoundedRangeOfIntegers}}
\subsection{Syntax}
\begin{flalign*}
\Npattern \derives\ & \Tleq \parsesep \Nexpr &
\end{flalign*}

\subsection{Abstract Syntax}
\RenderTypes[remove_hypertargets]{pattern_leq}

\ASTRuleDef{PLeq}
\begin{mathpar}
\inferrule{}{
  \buildpattern(\Npattern(\Tleq, \punnode{\Nexpr})) \astarrow
  \overname{\PatternLeq(\astof{\vexpr})}{\vastnode}
}
\end{mathpar}

\subsection{Typing}
\TypingRuleDef{PLeq}
\ExampleDef{Typing Less-or-equal Patterns}
\listingref{typing-pleq} shows examples of well-typed less-or-equal patterns.
\ASLListing{Well-typed less-or-equal patterns}{typing-pleq}{\typingtests/TypingRule.PLeq.asl}

\listingref{typing-pleq-bad} shows examples of ill-typed less-or-equal patterns.
\ASLListing{Ill-typed less-or-equal patterns}{typing-pleq-bad}{\typingtests/TypingRule.PLeq.bad.asl}


\RenderProseAndFormally{annotate_pattern_leq}
\CodeSubsection{\PLeqBegin}{\PLeqEnd}{../Typing.ml}

\subsection{Semantics}
\SemanticsRuleDef{PLeq}
\ExampleDef{Evaluation of a Less-or-equal Pattern}
In \listingref{semantics-pleq},
\texttt{match\_true} evaluates to \True, since \texttt{3} is less than or equal to \texttt{42},
whereas \texttt{match\_false} evaluates to \False, since \texttt{42} is not less than or equal to \texttt{3}.
\ASLListing{Matching against an upper bound value}{semantics-pleq}{\semanticstests/SemanticsRule.PLeq.asl}


\RenderProseAndFormally{eval_pattern_PLeq}
\CodeSubsection{\EvalPLeqBegin}{\EvalPLeqEnd}{../Interpreter.ml}

\section{Matching a Lower Bounded Range of Integers\label{sec:MatchingALowerBoundedRangeOfIntegers}}
\subsection{Syntax}
\begin{flalign*}
\Npattern \derives\ & \Tgeq \parsesep \Nexpr &
\end{flalign*}

\subsection{Abstract Syntax}
\RenderTypes[remove_hypertargets]{pattern_geq}

\ASTRuleDef{PGeq}
\begin{mathpar}
\inferrule{}{
  \buildpattern(\Npattern(\Tgeq, \punnode{\Nexpr})) \astarrow
  \overname{\PatternGeq(\astof{\vexpr})}{\vastnode}
}
\end{mathpar}

\subsection{Typing}
\TypingRuleDef{PGeq}
\ExampleDef{Typing Greater-or-equal Patterns}
\listingref{typing-pgeq} shows examples of well-typed greater-or-equal patterns.
\ASLListing{Well-typed greater-or-equal patterns}{typing-pgeq}{\typingtests/TypingRule.PGeq.asl}

\listingref{typing-pgeq-bad} shows examples of ill-typed greater-or-equal patterns.
\ASLListing{Ill-typed greater-or-equal patterns}{typing-pgeq-bad}{\typingtests/TypingRule.PGeq.bad.asl}


\RenderProseAndFormally{annotate_pattern_geq}
\CodeSubsection{\PGeqBegin}{\PGeqEnd}{../Typing.ml}

\subsection{Semantics}
\SemanticsRuleDef{PGeq}
\ExampleDef{Evaluation of a Greater-or-equal Pattern}
In \listingref{semantics-pgeq},
\texttt{match\_true} evaluates to \True, since \texttt{42} is greater than or equal to \texttt{3},
whereas \texttt{match\_false} evaluates to \False, since \texttt{3} is not greater than or equal to \texttt{42}.
\ASLListing{Matching against a lower bound}{semantics-pgeq}{\semanticstests/SemanticsRule.PGeq.asl}


\RenderProseAndFormally{eval_pattern_PGeq}
\CodeSubsection{\EvalPGeqBegin}{\EvalPGeqEnd}{../Interpreter.ml}

\section{Matching a Bitmask\label{sec:MatchingABitmask}}
Bitmasks are used to match sets of bitvectors where occurrences of $\xbit$,
and bit sequences in parentheses represent, don't-care values.

\subsection{Syntax}
\begin{flalign*}
\Npattern \derives\ & \Tmasklit &
\end{flalign*}

\subsection{Abstract Syntax}
\RenderTypes[remove_hypertargets]{pattern_mask}

\ASTRuleDef{PMask}
\begin{mathpar}
\inferrule{}{
  \buildpattern(\Npattern(\Tmasklit(\vm))) \astarrow
  \overname{\PatternMask(\vm)}{\vastnode}
}
\end{mathpar}

\subsection{Typing}
\TypingRuleDef{PMask}
\ExampleDef{Typing Mask Patterns}
\listingref{typing-pmask} shows examples of well-typed mask patterns.
\ASLListing{Well-typed mask patterns}{typing-pmask}{\typingtests/TypingRule.PMask.asl}

\listingref{typing-pmask-bad} shows examples of ill-typed mask patterns.
\ASLListing{An ill-typed mask pattern}{typing-pmask-bad}{\typingtests/TypingRule.PMask.bad.asl}


\RenderProseAndFormally{annotate_pattern_mask}
\CodeSubsection{\PMaskBegin}{\PMaskEnd}{../Typing.ml}


\subsection{Semantics}
\SemanticsRuleDef{PMask}
\ExampleDef{Evaluation of a Mask Pattern}
In \listingref{semantics-pmask}, \texttt{match\_true} evaluates to \True, since
\texttt{101010} matches the bitmask \texttt{xx1010},
whereas \texttt{match\_false} evaluates to \False, since
\texttt{101010} does not match the bitmask \texttt{0x1010}.
\ASLListing{Matching against a bitmask}{semantics-pmask}{\semanticstests/SemanticsRule.PMask.asl}

\RenderRelation{mask_match}
It is defined by the following table where the rows correspond to values of the bit and the columns correspond to mask values:
\[
\begin{array}{|c|c|c|c|}
  \hline
  \maskmatch & \zerobit & \onebit & \xbit\\
  \hline
  0 & \True & \False & \True\\
  \hline
  1 & \False & \True & \True\\
  \hline
\end{array}
\]

\RenderProseAndFormally{eval_pattern_PMask}
\CodeSubsection{\EvalPMaskBegin}{\EvalPMaskEnd}{../Interpreter.ml}


\section{Matching Pattern Lists\label{sec:MatchingPatternLists}}
Pattern-list forms are the only set-like pattern syntax. They appear at the top level
of an \texttt{IN} expression or a case alternative, and their elements are simple
patterns, that is, patterns derived from \Npattern. The forms \verb|{...}| and
\verb|!{...}| cannot be nested inside a pattern list.

\subsection{Syntax}
\begin{flalign*}
\Npatternset  \derives \  & \Tbnot \parsesep \Tlbrace \parsesep \Npatternlist \parsesep \Trbrace &\\
                  |\      & \Tlbrace \parsesep \Npatternlist \parsesep \Trbrace &\\
\Npatternlist \derives \ & \ClistOne{\Npattern} &
\end{flalign*}

\subsection{Abstract Syntax}
\RenderTypes{pattern_kind}

\ASTRuleDef{PatternSet}
\NotImportedToASLSpecYet{
\hypertarget{build-patternset}{}
The function
\[
  \buildpatternset(\overname{\parsenode{\Npatternset}}{\vparsednode}) \;\aslto\; \overname{(\overname{\pattern^*}{\vpatterns}, \overname{\patternkind}{\vpatternkind})}{\vastnode}
\]
transforms a parse node $\vparsednode$ into a pattern list and pattern kind.
} % END_OF_UNIMPORTED_RELATION

\begin{mathpar}
\inferrule[not]{
  \buildpatternlist(\punnode{\Npatternlist}) \astarrow \vpatternasts
}{
  \buildpatternset(\Npatternset(\Tbnot, \Tlbrace, \punnode{\Npatternlist}, \Trbrace)) \astarrow
  \overname{(\vpatternasts, \Negative)}{\vastnode}
}
\end{mathpar}

\begin{mathpar}
\inferrule[list]{
  \buildpatternlist(\punnode{\Npatternlist}) \astarrow \vpatternasts
}{
  \buildpatternset(\Npatternset(\Tlbrace, \punnode{\Npatternlist}, \Trbrace)) \astarrow
  \overname{(\vpatternasts, \Positive)}{\vastnode}
}
\end{mathpar}

\ASTRuleDef{PatternList}
\NotImportedToASLSpecYet{
\hypertarget{build-patternlist}{}
The function
\[
  \buildpatternlist(\overname{\parsenode{\Npatternlist}}{\vparsednode}) \;\aslto\; \overname{\pattern^*}{\vastnode}
\]
transforms a parse node $\vparsednode$ into a list of simple pattern AST nodes.
} % END_OF_UNIMPORTED_RELATION

\begin{mathpar}
\inferrule{
  \buildclist[\buildpattern](\vpatterns) \astarrow \vpatternasts
}{
  \buildpatternlist(\Npatternlist(\namednode{\vpatterns}{\ClistOne{\Npattern}})) \astarrow
  \overname{\vpatternasts}{\vastnode}
}
\end{mathpar}

\subsection{Typing}
\RenderRelation{annotate_pattern_list_and_kind}
\RenderProseAndFormally{annotate_pattern_list_and_kind}

\subsection{Semantics}
\SemanticsRuleDef{EvalPatternListAndKind}
\ExampleDef{Evaluation of Pattern Lists}

A positive pattern list matches when at least one pattern in the list matches the
discriminant value. A negative pattern list matches when no pattern in the list
matches the discriminant value.

\RenderRelation{eval_pattern_list_and_kind}
\RenderProseAndFormally{eval_pattern_list_and_kind}
\CodeSubsection{\EvalPatternListAndKindBegin}{\EvalPatternListAndKindEnd}{../Interpreter.ml}
